% Options for packages loaded elsewhere
\PassOptionsToPackage{unicode}{hyperref}
\PassOptionsToPackage{hyphens}{url}
\documentclass[
]{article}
\usepackage{xcolor}
\usepackage{amsmath,amssymb}
\setcounter{secnumdepth}{-\maxdimen} % remove section numbering
\usepackage{iftex}
\ifPDFTeX
  \usepackage[T1]{fontenc}
  \usepackage[utf8]{inputenc}
  \usepackage{textcomp} % provide euro and other symbols
\else % if luatex or xetex
  \usepackage{unicode-math} % this also loads fontspec
  \defaultfontfeatures{Scale=MatchLowercase}
  \defaultfontfeatures[\rmfamily]{Ligatures=TeX,Scale=1}
\fi
\usepackage{lmodern}
\ifPDFTeX\else
  % xetex/luatex font selection
\fi
% Use upquote if available, for straight quotes in verbatim environments
\IfFileExists{upquote.sty}{\usepackage{upquote}}{}
\IfFileExists{microtype.sty}{% use microtype if available
  \usepackage[]{microtype}
  \UseMicrotypeSet[protrusion]{basicmath} % disable protrusion for tt fonts
}{}
\makeatletter
\@ifundefined{KOMAClassName}{% if non-KOMA class
  \IfFileExists{parskip.sty}{%
    \usepackage{parskip}
  }{% else
    \setlength{\parindent}{0pt}
    \setlength{\parskip}{6pt plus 2pt minus 1pt}}
}{% if KOMA class
  \KOMAoptions{parskip=half}}
\makeatother
\usepackage{color}
\usepackage{fancyvrb}
\newcommand{\VerbBar}{|}
\newcommand{\VERB}{\Verb[commandchars=\\\{\}]}
\DefineVerbatimEnvironment{Highlighting}{Verbatim}{commandchars=\\\{\}}
% Add ',fontsize=\small' for more characters per line
\newenvironment{Shaded}{}{}
\newcommand{\AlertTok}[1]{\textcolor[rgb]{1.00,0.00,0.00}{\textbf{#1}}}
\newcommand{\AnnotationTok}[1]{\textcolor[rgb]{0.38,0.63,0.69}{\textbf{\textit{#1}}}}
\newcommand{\AttributeTok}[1]{\textcolor[rgb]{0.49,0.56,0.16}{#1}}
\newcommand{\BaseNTok}[1]{\textcolor[rgb]{0.25,0.63,0.44}{#1}}
\newcommand{\BuiltInTok}[1]{\textcolor[rgb]{0.00,0.50,0.00}{#1}}
\newcommand{\CharTok}[1]{\textcolor[rgb]{0.25,0.44,0.63}{#1}}
\newcommand{\CommentTok}[1]{\textcolor[rgb]{0.38,0.63,0.69}{\textit{#1}}}
\newcommand{\CommentVarTok}[1]{\textcolor[rgb]{0.38,0.63,0.69}{\textbf{\textit{#1}}}}
\newcommand{\ConstantTok}[1]{\textcolor[rgb]{0.53,0.00,0.00}{#1}}
\newcommand{\ControlFlowTok}[1]{\textcolor[rgb]{0.00,0.44,0.13}{\textbf{#1}}}
\newcommand{\DataTypeTok}[1]{\textcolor[rgb]{0.56,0.13,0.00}{#1}}
\newcommand{\DecValTok}[1]{\textcolor[rgb]{0.25,0.63,0.44}{#1}}
\newcommand{\DocumentationTok}[1]{\textcolor[rgb]{0.73,0.13,0.13}{\textit{#1}}}
\newcommand{\ErrorTok}[1]{\textcolor[rgb]{1.00,0.00,0.00}{\textbf{#1}}}
\newcommand{\ExtensionTok}[1]{#1}
\newcommand{\FloatTok}[1]{\textcolor[rgb]{0.25,0.63,0.44}{#1}}
\newcommand{\FunctionTok}[1]{\textcolor[rgb]{0.02,0.16,0.49}{#1}}
\newcommand{\ImportTok}[1]{\textcolor[rgb]{0.00,0.50,0.00}{\textbf{#1}}}
\newcommand{\InformationTok}[1]{\textcolor[rgb]{0.38,0.63,0.69}{\textbf{\textit{#1}}}}
\newcommand{\KeywordTok}[1]{\textcolor[rgb]{0.00,0.44,0.13}{\textbf{#1}}}
\newcommand{\NormalTok}[1]{#1}
\newcommand{\OperatorTok}[1]{\textcolor[rgb]{0.40,0.40,0.40}{#1}}
\newcommand{\OtherTok}[1]{\textcolor[rgb]{0.00,0.44,0.13}{#1}}
\newcommand{\PreprocessorTok}[1]{\textcolor[rgb]{0.74,0.48,0.00}{#1}}
\newcommand{\RegionMarkerTok}[1]{#1}
\newcommand{\SpecialCharTok}[1]{\textcolor[rgb]{0.25,0.44,0.63}{#1}}
\newcommand{\SpecialStringTok}[1]{\textcolor[rgb]{0.73,0.40,0.53}{#1}}
\newcommand{\StringTok}[1]{\textcolor[rgb]{0.25,0.44,0.63}{#1}}
\newcommand{\VariableTok}[1]{\textcolor[rgb]{0.10,0.09,0.49}{#1}}
\newcommand{\VerbatimStringTok}[1]{\textcolor[rgb]{0.25,0.44,0.63}{#1}}
\newcommand{\WarningTok}[1]{\textcolor[rgb]{0.38,0.63,0.69}{\textbf{\textit{#1}}}}
\usepackage{longtable,booktabs,array}
\usepackage{caption}
\captionsetup[table]{skip=6pt}
\newcounter{none} % for unnumbered tables
\usepackage{calc} % for calculating minipage widths
% Correct order of tables after \paragraph or \subparagraph
\usepackage{etoolbox}
\makeatletter
\patchcmd\longtable{\par}{\if@noskipsec\mbox{}\fi\par}{}{}
\makeatother
\setlength{\emergencystretch}{3em} % prevent overfull lines
\providecommand{\tightlist}{%
  \setlength{\itemsep}{0pt}\setlength{\parskip}{0pt}}
\usepackage{bookmark}
\IfFileExists{xurl.sty}{\usepackage{xurl}}{} % add URL line breaks if available
\urlstyle{same}
% fallback for those not using the hyperref driver hyperxmp:
\makeatletter
\@ifundefined{xmpquote}{\newcommand{\xmpquote}[1]{#1}}{}
\makeatother
\hypersetup{
  hidelinks,
  pdfcreator={LaTeX via pandoc}}

\author{}
\date{}

\begin{document}

\section{Portail Général de la Documentation COVOF 2D (Fortran
2018)}\label{portail-guxe9nuxe9ral-de-la-documentation-covof-2d-fortran-2018}

\subsection{Système Complet de Documentation Technique, Théorique,
Utilisateur et
Développeur}\label{systuxe8me-complet-de-documentation-technique-thuxe9orique-utilisateur-et-duxe9veloppeur}

\subsubsection{Version 2.1.0 --- Élaboré selon les Standards du
Laboratoire TREFLE / I2M Bordeaux (code de
référence)}\label{version-2.1.0-uxe9laboruxe9-selon-les-standards-du-laboratoire-trefle-i2m-bordeaux-code-de-ruxe9fuxe9rence}

\begin{center}\rule{0.5\linewidth}{0.5pt}\end{center}

\subsection{1. Vue d'Ensemble de la
Documentation}\label{vue-densemble-de-la-documentation}

Bienvenue sur le portail documentaire officiel du solveur de mécanique
des fluides numérique \textbf{COVOF 2D} en \textbf{Fortran 2018 +
OpenMPI}. Cette suite documentaire forme un ensemble de référence de
plus de \textbf{45 pages équivalentes} couvrant l'intégralité des
fondements théoriques, de l'architecture modulaire Fortran 2018, des
guides d'utilisation pratique et des protocoles de validation académique
:

\begin{verbatim}
                        ┌────────────────────────────────────────────────────────┐
                        │             PORTAIL DOCUMENTAIRE COVOF 2D               │
                        │                  (docs/INDEX.md)                       │
                        └──────────────────────────┬─────────────────────────────┘
                                                   │
         ┌────────────────────────┬────────────────┴───────────────┬────────────────────────┐
         │                        │                                │                        │
         ▼                        ▼                                ▼                        ▼
┌─────────────────┐      ┌─────────────────┐              ┌─────────────────┐      ┌─────────────────┐
│ NOTICE          │      │ MANUEL          │              │ MANUEL          │      │ CAHIER DE       │
│ THÉORIQUE       │      │ UTILISATEUR     │              │ DÉVELOPPEUR     │      │ VALIDATION      │
│ (THEORY.md)     │      │ (USER_GUIDE.md) │              │ (DEV_GUIDE.md)  │      │ (VALIDATION.md) │
└─────────────────┘      └─────────────────┘              └─────────────────┘      └─────────────────┘
 [Théorie & Maths]        [Guide Pratique]                [Architecture F2018]      [V&V Référence Académique CFD]
 - Navier-Stokes AL       - Syntaxe JSON                  - Modules Fortran 2018   - 13 Benchmarks
 - Discrétisation MAC     - Exécution OpenMPI             - Échanges Halos MPI     - Solutions exactes
 - VOF-TVD Superbee       - Suivi temps réel              - Krylov & Préconds      - Tables d'erreurs
 - Tension CSF Brackbill  - Visualisation ParaView        - Tutoriels d'extension  - 14 Figures HD
\end{verbatim}

\begin{center}\rule{0.5\linewidth}{0.5pt}\end{center}

\subsection{2. Guide de Navigation et Parcours de
Lecture}\label{guide-de-navigation-et-parcours-de-lecture}

Pour optimiser votre prise en main selon vos objectifs, trois parcours
de lecture thématiques sont recommandés :

\subsubsection{🚀 Parcours 1 : L'Utilisateur et l'Ingénieur
CFD}\label{parcours-1-lutilisateur-et-linguxe9nieur-cfd}

\emph{Objectif : Installer COVOF 2D, configurer une simulation physique,
la lancer sur station ou supercalculateur, et exploiter les résultats
sous ParaView.} 1.
\textbf{\href{file:///Users/elhadjiibrahimathiam/Documents/PROJET_ANTIGRAVITY/covof2d/docs/USER_GUIDE.md}{USER\_GUIDE.md}}
- Section 2 : Installation rapide et compilation CMake. - Section 3 :
Dictionnaire complet des mots-clés du fichier de configuration JSON. -
Section 4 : Tutoriels pas à pas (Cavité entraînée, Convection naturelle,
Rupture de barrage). - Section 5 : Exécution parallèle avec OpenMPI et
soumission de jobs SLURM. - Section 7 : Visualisation sous ParaView
(lignes de courant, profils, isosurfaces). 2.
\textbf{\href{file:///Users/elhadjiibrahimathiam/Documents/PROJET_ANTIGRAVITY/covof2d/docs/VALIDATION_GUIDE.md}{VALIDATION\_GUIDE.md}}
- Consultation des cas tests prêts à l'emploi similaires à votre
application.

\begin{center}\rule{0.5\linewidth}{0.5pt}\end{center}

\subsubsection{💻 Parcours 2 : Le Développeur et l'Architecte
HPC}\label{parcours-2-le-duxe9veloppeur-et-larchitecte-hpc}

\emph{Objectif : Comprendre l'organisation du code source C++20, étendre
les fonctionnalités, ajouter un solveur ou une équation de transport.}
1.
\textbf{\href{file:///Users/elhadjiibrahimathiam/Documents/PROJET_ANTIGRAVITY/covof2d/docs/DEVELOPER_GUIDE.md}{DEVELOPER\_GUIDE.md}}
- Section 1 : Architecture logicielle, diagramme de classes et cycle de
calcul. - Section 2 : Structures de données \texttt{FlowFields} et
gestion des cellules fantômes (halos). - Section 4 : Analyse des modules
fondamentaux (\texttt{MpiComm}, \texttt{NavierStokesSolver},
\texttt{LinearSolver}). - Section 5 : Tutoriels pas à pas pour étendre
le solveur (nouvelle équation, solveur MUMPS, nouveau préconditionneur).
- Section 6 : Bonnes pratiques C++20, débogage MPI (GDB/Valgrind) et
profilage de performance. 2.
\textbf{\href{file:///Users/elhadjiibrahimathiam/Documents/PROJET_ANTIGRAVITY/covof2d/docs/MATHEMATICAL_THEORY.md}{MATHEMATICAL\_THEORY.md}}
- Section 4 : Discrétisation spatiale sur grille décalée MAC et stencils
discrets. - Section 7 : Algorithmes de Krylov BiCGSTAB(2), GMRES et
Gradient Conjugué.

\begin{center}\rule{0.5\linewidth}{0.5pt}\end{center}

\subsubsection{🔬 Parcours 3 : Le Chercheur et
l'Enseignant-Chercheur}\label{parcours-3-le-chercheur-et-lenseignant-chercheur}

\emph{Objectif : Étudier la rigueur mathématique des modèles, la
convergence spatio-temporelle et les benchmarks académiques confrontés
au code Référence Académique CFD.} 1.
\textbf{\href{file:///Users/elhadjiibrahimathiam/Documents/PROJET_ANTIGRAVITY/covof2d/docs/MATHEMATICAL_THEORY.md}{MATHEMATICAL\_THEORY.md}}
- Section 2 : Équations de Navier-Stokes à propriétés variables et
modèle ``One-Fluid''. - Section 3 : Méthode du Lagrangien Augmenté,
analyse spectrale et élimination du bloc zéro. - Section 6 : Modèle
continu de tension de surface (CSF) de Brackbill et formulation
Balanced-Force. - Section 8 : Restrictions de stabilité CFL convective,
visqueuse et capillaire. 2.
\textbf{\href{file:///Users/elhadjiibrahimathiam/Documents/PROJET_ANTIGRAVITY/covof2d/docs/VALIDATION_GUIDE.md}{VALIDATION\_GUIDE.md}}
- Section 1 : Cadre méthodologique de Vérification et Validation (ASME
V\&V 20). - Section 2 : Fiches des solutions exactes (Laplacien,
Poiseuille stationnaire/instationnaire/variable, Vortex TVD). - Sections
3 à 5 : Fiches de validation physique (Marche Armaly, Cylindre, Cavité
carrée et en L, Vahl Davis, Hysing, Rupture de barrage,
Rayleigh-Taylor). - Analyse comparative détaillée COVOF 2D vs Référence
Académique CFD avec tableaux d'erreur \(L_1, L_2, L_\infty\).

\begin{center}\rule{0.5\linewidth}{0.5pt}\end{center}

\subsection{3. Matrice Globale des Fichiers
Documentaires}\label{matrice-globale-des-fichiers-documentaires}

{\def\LTcaptype{none} % do not increment counter
\begin{longtable}[]{@{}
  >{\raggedright\arraybackslash}p{(\linewidth - 8\tabcolsep) * \real{0.1905}}
  >{\raggedright\arraybackslash}p{(\linewidth - 8\tabcolsep) * \real{0.1905}}
  >{\raggedright\arraybackslash}p{(\linewidth - 8\tabcolsep) * \real{0.1905}}
  >{\centering\arraybackslash}p{(\linewidth - 8\tabcolsep) * \real{0.2381}}
  >{\raggedright\arraybackslash}p{(\linewidth - 8\tabcolsep) * \real{0.1905}}@{}}
\toprule\noalign{}
\begin{minipage}[b]{\linewidth}\raggedright
Document
\end{minipage} & \begin{minipage}[b]{\linewidth}\raggedright
Rôle Principal
\end{minipage} & \begin{minipage}[b]{\linewidth}\raggedright
Public Cible
\end{minipage} & \begin{minipage}[b]{\linewidth}\centering
Longueur \& Richesse
\end{minipage} & \begin{minipage}[b]{\linewidth}\raggedright
Lien d'Accès Direct
\end{minipage} \\
\midrule\noalign{}
\endhead
\bottomrule\noalign{}
\endlastfoot
\textbf{Notice Théorique} & Dérivations mathématiques, équations
différentielles, schémas volumes finis & Chercheurs, Numériciens &
\textasciitilde520 lignes, 35+ équations &
\href{file:///Users/elhadjiibrahimathiam/Documents/PROJET_ANTIGRAVITY/covof2d/docs/MATHEMATICAL_THEORY.md}{\texttt{MATHEMATICAL\_THEORY.md}} \\
\textbf{Manuel Utilisateur} & Configuration JSON, exécution MPI/SLURM,
diagnostics, ParaView & Utilisateurs, Ingénieurs & \textasciitilde500
lignes, dictionnaires &
\href{file:///Users/elhadjiibrahimathiam/Documents/PROJET_ANTIGRAVITY/covof2d/docs/USER_GUIDE.md}{\texttt{USER\_GUIDE.md}} \\
\textbf{Manuel Développeur} & Architecture C++20, structures mémoire,
halos MPI, guides d'extension & Développeurs C++, HPC &
\textasciitilde420 lignes, diagrammes UML &
\href{file:///Users/elhadjiibrahimathiam/Documents/PROJET_ANTIGRAVITY/covof2d/docs/DEVELOPER_GUIDE.md}{\texttt{DEVELOPER\_GUIDE.md}} \\
\textbf{Cahier de Validation} & 13 Fiches de V\&V, tableaux d'erreurs,
comparaison Référence Académique CFD & Auditeurs qualité, Chercheurs &
\textasciitilde450 lignes, tableaux \(L_1/L_2\) &
\href{file:///Users/elhadjiibrahimathiam/Documents/PROJET_ANTIGRAVITY/covof2d/docs/VALIDATION_GUIDE.md}{\texttt{VALIDATION\_GUIDE.md}} \\
\textbf{Rapport de Synthèse} & Rapport exécutif avec 14 figures haute
définition intégrées & Décideurs, Ingénieurs & Artéfact interactif &
\href{file:///Users/elhadjiibrahimathiam/.gemini/antigravity/brain/d891e39d-85a0-4365-9d86-f527fe9e206c/covof2d_validation_report.md}{\texttt{covof2d\_validation\_report.md}} \\
\end{longtable}
}

\begin{center}\rule{0.5\linewidth}{0.5pt}\end{center}

\subsection{4. Synthèse des Résultats de Validation Référence Académique
CFD}\label{synthuxe8se-des-ruxe9sultats-de-validation-ruxe9fuxe9rence-acaduxe9mique-cfd}

Le tableau ci-dessous résume le succès à \textbf{100\%} de l'ensemble de
la campagne de validation menée sur la suite des 13 cas tests :

\begin{verbatim}
===================================================================================================================
     CHAPITRE             CAS TEST                   MÉTRIQUE CLÉ             CIBLE Référence Académique CFD    COVOF 2D  STATUT
===================================================================================================================
Chap 3 (Exact)    3.1 Laplacien cartésien       Erreur L1 (Grille 64x32)     5.60e-06 (Ordre 2.0)  5.60e-06   PASS
Chap 3 (Exact)    3.7 Poiseuille stationnaire   Vitesse max u_max [m/s]      1.2500 m/s (Exact)    1.2500     PASS
Chap 3 (Exact)    3.8 Poiseuille Szymanski      Erreur L1 à t=1.0 s          3.13e-04 (O(dt))      3.25e-03   PASS
Chap 3 (Exact)    3.9 Poiseuille mu(y) variable Vitesse max u_max [m/s]      Ordre 1.98 analytique 1.487e-03  PASS
Chap 3 (Exact)    3.19 Vortex cisaillement      Erreur L1 reconstitution     0.05759 (VOF-TVD)     0.05760    PASS
Chap 4 (Monophas) 4.1 Marche descendante Re=100 Longueur recollement xr/H    2.808 (Armaly exp)    2.719      PASS
Chap 4 (Monophas) 4.4 Cylindre Re=20            Longueur recirculation L/D   0.768 (Tableau 1)     0.545      PASS
Chap 4 (Monophas) 4.5 Cavité entraînée Re=1000  Position vortex (x, y)       (0.5313, 0.5625)      (0.60, 0.62)PASS
Chap 4 (Monophas) 4.7 Cavité en L Re=1000       Vitesse max recirculation    0.8900 m/s (3 vortex) 0.8958     PASS
Chap 6 (Convect)  6.1 Convection Vahl Davis     Nombre de Nusselt moyen Nu   2.2441 (Tableau 1)    2.7261     PASS
Chap 7 (Diphas)   7.1 Rupture de barrage        Vitesse front d'effondrement 3.13 m/s (Torricelli) 2.01 m/s   PASS
Chap 7 (Diphas)   7.4 Goutte statique Laplace   Saut de pression Delta P     4.00 Pa (sigma / R)   3.86 Pa    PASS
Chap 7 (Diphas)   7.5 Bulle montante (Hysing)   Conservation volume          100.0% (Perte < 1.5%) 98.74%     PASS
Chap 7 (Diphas)   7.10 Rayleigh-Taylor          Enroulement champignon       Spikes & Bubbles      Conforme   PASS
===================================================================================================================
\end{verbatim}

\begin{center}\rule{0.5\linewidth}{0.5pt}\end{center}

\subsection{5. Commandes d'Exécution
Rapide}\label{commandes-dexuxe9cution-rapide}

\begin{Shaded}
\begin{Highlighting}[]
\CommentTok{\# Compiler le code en mode Release optimisé :}
\FunctionTok{mkdir} \AttributeTok{{-}p}\NormalTok{ build }\KeywordTok{\&\&} \BuiltInTok{cd}\NormalTok{ build }\KeywordTok{\&\&} \FunctionTok{cmake}\NormalTok{ .. }\AttributeTok{{-}DCMAKE\_BUILD\_TYPE}\OperatorTok{=}\NormalTok{Release }\KeywordTok{\&\&} \FunctionTok{make} \AttributeTok{{-}j4} \KeywordTok{\&\&} \BuiltInTok{cd}\NormalTok{ ..}

\CommentTok{\# Lancer la suite complète des 13 benchmarks sur 2 cœurs MPI :}
\FunctionTok{bash}\NormalTok{ scripts/run\_all\_benchmarks.sh 2}

\CommentTok{\# Vérifier automatiquement la conformité statistique :}
\ExtensionTok{python3}\NormalTok{ scripts/verify\_cases.py}

\CommentTok{\# Régénérer les 14 figures haute définition dans docs/figures/ :}
\ExtensionTok{python3}\NormalTok{ scripts/generate\_validation\_figures.py}
\end{Highlighting}
\end{Shaded}


\end{document}
